% Macros shared by the test documents, read with % Macros shared by the test documents, read with % Macros shared by the test documents, read with % Macros shared by the test documents, read with \input{tests/checks}.
%
% A test document uses the class as any letter does. It puts \Mark{label}
% into the texts that the class typesets and compares the positions of the
% marks, known from the previous run, with the expected ones. Every check
% writes one line to the log, "TEST OK: ..." or "TEST FAILED: ...", and
% tests/run.sh looks for these lines after the last run. A failed check
% must not stop the compilation: the positions in the footer are correct
% only from the third run on.

\RequirePackage{zref-savepos}

\makeatletter

% \Mark{label} remembers the place where it is typeset. \PageMark{label}
% does the same with the number of the page appended to the label.
\DeclareRobustCommand{\Mark}[1]{\leavevmode\zsavepos{#1}}
\DeclareRobustCommand{\PageMark}[1]{\Mark{#1\arabic{page}}}
\pdfstringdefDisableCommands{\let\Mark\@gobble\let\PageMark\@gobble}

% The position of a mark, measured from the left and from the top edge of
% the page. Can be used in the arguments of \CheckPos.
\newcommand{\PosX}[1]{\dimexpr\zposx{#1}sp\relax}
\newcommand{\PosY}[1]{\dimexpr\paperheight-\zposy{#1}sp\relax}

% The family, the series and the size of the current font, for \CheckEqual.
\newcommand{\FontFamily}{\f@family}
\newcommand{\FontSeries}{\f@series}
\newcommand{\FontSize}{\f@size}

% A length in millimetres, for the messages.
\newcommand{\Check@mm}[1]{\strip@pt\dimexpr(#1)*2540/7227\relax}

% \CheckPos[tolerance]{label}{x}{y} checks that the mark is x from the left
% edge and y from the top edge of the page.
\newcommand{\CheckPos}[4][0.02mm]{%
	\@tempswatrue
	\@tempdima=\dimexpr\PosX{#2}-(#3)\relax
	\ifdim\@tempdima<\z@ \@tempdima=-\@tempdima \fi
	\ifdim\@tempdima>\dimexpr#1\relax \@tempswafalse \fi
	\@tempdima=\dimexpr\PosY{#2}-(#4)\relax
	\ifdim\@tempdima<\z@ \@tempdima=-\@tempdima \fi
	\ifdim\@tempdima>\dimexpr#1\relax \@tempswafalse \fi
	\if@tempswa
		\typeout{TEST OK: position of #2}%
	\else
		\typeout{TEST FAILED: position of #2:
			expected (\Check@mm{#3}, \Check@mm{#4}) mm,
			got (\Check@mm{\PosX{#2}}, \Check@mm{\PosY{#2}}) mm}%
	\fi
}

% \CheckAbsent{label} checks that the mark was not typeset at all.
\newcommand{\CheckAbsent}[1]{%
	\zref@ifrefundefined{#1}%
		{\typeout{TEST OK: no #1}}%
		{\typeout{TEST FAILED: #1 should be absent, it is at
			(\Check@mm{\PosX{#1}}, \Check@mm{\PosY{#1}}) mm}}%
}

% \CheckEqual{what}{actual}{expected} compares two texts after expansion.
\newcommand{\CheckEqual}[3]{%
	\protected@edef\Check@a{#2}%
	\protected@edef\Check@b{#3}%
	\ifx\Check@a\Check@b
		\typeout{TEST OK: #1}%
	\else
		\typeout{TEST FAILED: #1: expected "\Check@b", got "\Check@a"}%
	\fi
}

% \CheckColour{name}{R,G,B} checks the definition of a colour; the name .
% stands for the colour of the text at the place of the check.
\newcommand{\CheckColour}[2]{%
	\extractcolorspecs{#1}{\Check@a}{\Check@b}%
	\convertcolorspec{\Check@a}{\Check@b}{RGB}{\Check@c}%
	\CheckEqual{colour #1}{\Check@c}{#2}%
}

\makeatother
.
%
% A test document uses the class as any letter does. It puts \Mark{label}
% into the texts that the class typesets and compares the positions of the
% marks, known from the previous run, with the expected ones. Every check
% writes one line to the log, "TEST OK: ..." or "TEST FAILED: ...", and
% tests/run.sh looks for these lines after the last run. A failed check
% must not stop the compilation: the positions in the footer are correct
% only from the third run on.

\RequirePackage{zref-savepos}

\makeatletter

% \Mark{label} remembers the place where it is typeset. \PageMark{label}
% does the same with the number of the page appended to the label.
\DeclareRobustCommand{\Mark}[1]{\leavevmode\zsavepos{#1}}
\DeclareRobustCommand{\PageMark}[1]{\Mark{#1\arabic{page}}}
\pdfstringdefDisableCommands{\let\Mark\@gobble\let\PageMark\@gobble}

% The position of a mark, measured from the left and from the top edge of
% the page. Can be used in the arguments of \CheckPos.
\newcommand{\PosX}[1]{\dimexpr\zposx{#1}sp\relax}
\newcommand{\PosY}[1]{\dimexpr\paperheight-\zposy{#1}sp\relax}

% The family, the series and the size of the current font, for \CheckEqual.
\newcommand{\FontFamily}{\f@family}
\newcommand{\FontSeries}{\f@series}
\newcommand{\FontSize}{\f@size}

% A length in millimetres, for the messages.
\newcommand{\Check@mm}[1]{\strip@pt\dimexpr(#1)*2540/7227\relax}

% \CheckPos[tolerance]{label}{x}{y} checks that the mark is x from the left
% edge and y from the top edge of the page.
\newcommand{\CheckPos}[4][0.02mm]{%
	\@tempswatrue
	\@tempdima=\dimexpr\PosX{#2}-(#3)\relax
	\ifdim\@tempdima<\z@ \@tempdima=-\@tempdima \fi
	\ifdim\@tempdima>\dimexpr#1\relax \@tempswafalse \fi
	\@tempdima=\dimexpr\PosY{#2}-(#4)\relax
	\ifdim\@tempdima<\z@ \@tempdima=-\@tempdima \fi
	\ifdim\@tempdima>\dimexpr#1\relax \@tempswafalse \fi
	\if@tempswa
		\typeout{TEST OK: position of #2}%
	\else
		\typeout{TEST FAILED: position of #2:
			expected (\Check@mm{#3}, \Check@mm{#4}) mm,
			got (\Check@mm{\PosX{#2}}, \Check@mm{\PosY{#2}}) mm}%
	\fi
}

% \CheckAbsent{label} checks that the mark was not typeset at all.
\newcommand{\CheckAbsent}[1]{%
	\zref@ifrefundefined{#1}%
		{\typeout{TEST OK: no #1}}%
		{\typeout{TEST FAILED: #1 should be absent, it is at
			(\Check@mm{\PosX{#1}}, \Check@mm{\PosY{#1}}) mm}}%
}

% \CheckEqual{what}{actual}{expected} compares two texts after expansion.
\newcommand{\CheckEqual}[3]{%
	\protected@edef\Check@a{#2}%
	\protected@edef\Check@b{#3}%
	\ifx\Check@a\Check@b
		\typeout{TEST OK: #1}%
	\else
		\typeout{TEST FAILED: #1: expected "\Check@b", got "\Check@a"}%
	\fi
}

% \CheckColour{name}{R,G,B} checks the definition of a colour; the name .
% stands for the colour of the text at the place of the check.
\newcommand{\CheckColour}[2]{%
	\extractcolorspecs{#1}{\Check@a}{\Check@b}%
	\convertcolorspec{\Check@a}{\Check@b}{RGB}{\Check@c}%
	\CheckEqual{colour #1}{\Check@c}{#2}%
}

\makeatother
.
%
% A test document uses the class as any letter does. It puts \Mark{label}
% into the texts that the class typesets and compares the positions of the
% marks, known from the previous run, with the expected ones. Every check
% writes one line to the log, "TEST OK: ..." or "TEST FAILED: ...", and
% tests/run.sh looks for these lines after the last run. A failed check
% must not stop the compilation: the positions in the footer are correct
% only from the third run on.

\RequirePackage{zref-savepos}

\makeatletter

% \Mark{label} remembers the place where it is typeset. \PageMark{label}
% does the same with the number of the page appended to the label.
\DeclareRobustCommand{\Mark}[1]{\leavevmode\zsavepos{#1}}
\DeclareRobustCommand{\PageMark}[1]{\Mark{#1\arabic{page}}}
\pdfstringdefDisableCommands{\let\Mark\@gobble\let\PageMark\@gobble}

% The position of a mark, measured from the left and from the top edge of
% the page. Can be used in the arguments of \CheckPos.
\newcommand{\PosX}[1]{\dimexpr\zposx{#1}sp\relax}
\newcommand{\PosY}[1]{\dimexpr\paperheight-\zposy{#1}sp\relax}

% The family, the series and the size of the current font, for \CheckEqual.
\newcommand{\FontFamily}{\f@family}
\newcommand{\FontSeries}{\f@series}
\newcommand{\FontSize}{\f@size}

% A length in millimetres, for the messages.
\newcommand{\Check@mm}[1]{\strip@pt\dimexpr(#1)*2540/7227\relax}

% \CheckPos[tolerance]{label}{x}{y} checks that the mark is x from the left
% edge and y from the top edge of the page.
\newcommand{\CheckPos}[4][0.02mm]{%
	\@tempswatrue
	\@tempdima=\dimexpr\PosX{#2}-(#3)\relax
	\ifdim\@tempdima<\z@ \@tempdima=-\@tempdima \fi
	\ifdim\@tempdima>\dimexpr#1\relax \@tempswafalse \fi
	\@tempdima=\dimexpr\PosY{#2}-(#4)\relax
	\ifdim\@tempdima<\z@ \@tempdima=-\@tempdima \fi
	\ifdim\@tempdima>\dimexpr#1\relax \@tempswafalse \fi
	\if@tempswa
		\typeout{TEST OK: position of #2}%
	\else
		\typeout{TEST FAILED: position of #2:
			expected (\Check@mm{#3}, \Check@mm{#4}) mm,
			got (\Check@mm{\PosX{#2}}, \Check@mm{\PosY{#2}}) mm}%
	\fi
}

% \CheckAbsent{label} checks that the mark was not typeset at all.
\newcommand{\CheckAbsent}[1]{%
	\zref@ifrefundefined{#1}%
		{\typeout{TEST OK: no #1}}%
		{\typeout{TEST FAILED: #1 should be absent, it is at
			(\Check@mm{\PosX{#1}}, \Check@mm{\PosY{#1}}) mm}}%
}

% \CheckEqual{what}{actual}{expected} compares two texts after expansion.
\newcommand{\CheckEqual}[3]{%
	\protected@edef\Check@a{#2}%
	\protected@edef\Check@b{#3}%
	\ifx\Check@a\Check@b
		\typeout{TEST OK: #1}%
	\else
		\typeout{TEST FAILED: #1: expected "\Check@b", got "\Check@a"}%
	\fi
}

% \CheckColour{name}{R,G,B} checks the definition of a colour; the name .
% stands for the colour of the text at the place of the check.
\newcommand{\CheckColour}[2]{%
	\extractcolorspecs{#1}{\Check@a}{\Check@b}%
	\convertcolorspec{\Check@a}{\Check@b}{RGB}{\Check@c}%
	\CheckEqual{colour #1}{\Check@c}{#2}%
}

\makeatother
.
%
% A test document uses the class as any letter does. It puts \Mark{label}
% into the texts that the class typesets and compares the positions of the
% marks, known from the previous run, with the expected ones. Every check
% writes one line to the log, "TEST OK: ..." or "TEST FAILED: ...", and
% tests/run.sh looks for these lines after the last run. A failed check
% must not stop the compilation: the positions in the footer are correct
% only from the third run on.

\RequirePackage{zref-savepos}

\makeatletter

% \Mark{label} remembers the place where it is typeset. \PageMark{label}
% does the same with the number of the page appended to the label.
\DeclareRobustCommand{\Mark}[1]{\leavevmode\zsavepos{#1}}
\DeclareRobustCommand{\PageMark}[1]{\Mark{#1\arabic{page}}}
\pdfstringdefDisableCommands{\let\Mark\@gobble\let\PageMark\@gobble}

% The position of a mark, measured from the left and from the top edge of
% the page. Can be used in the arguments of \CheckPos.
\newcommand{\PosX}[1]{\dimexpr\zposx{#1}sp\relax}
\newcommand{\PosY}[1]{\dimexpr\paperheight-\zposy{#1}sp\relax}

% The family, the series and the size of the current font, for \CheckEqual.
\newcommand{\FontFamily}{\f@family}
\newcommand{\FontSeries}{\f@series}
\newcommand{\FontSize}{\f@size}

% A length in millimetres, for the messages.
\newcommand{\Check@mm}[1]{\strip@pt\dimexpr(#1)*2540/7227\relax}

% \CheckPos[tolerance]{label}{x}{y} checks that the mark is x from the left
% edge and y from the top edge of the page.
\newcommand{\CheckPos}[4][0.02mm]{%
	\@tempswatrue
	\@tempdima=\dimexpr\PosX{#2}-(#3)\relax
	\ifdim\@tempdima<\z@ \@tempdima=-\@tempdima \fi
	\ifdim\@tempdima>\dimexpr#1\relax \@tempswafalse \fi
	\@tempdima=\dimexpr\PosY{#2}-(#4)\relax
	\ifdim\@tempdima<\z@ \@tempdima=-\@tempdima \fi
	\ifdim\@tempdima>\dimexpr#1\relax \@tempswafalse \fi
	\if@tempswa
		\typeout{TEST OK: position of #2}%
	\else
		\typeout{TEST FAILED: position of #2:
			expected (\Check@mm{#3}, \Check@mm{#4}) mm,
			got (\Check@mm{\PosX{#2}}, \Check@mm{\PosY{#2}}) mm}%
	\fi
}

% \CheckAbsent{label} checks that the mark was not typeset at all.
\newcommand{\CheckAbsent}[1]{%
	\zref@ifrefundefined{#1}%
		{\typeout{TEST OK: no #1}}%
		{\typeout{TEST FAILED: #1 should be absent, it is at
			(\Check@mm{\PosX{#1}}, \Check@mm{\PosY{#1}}) mm}}%
}

% \CheckEqual{what}{actual}{expected} compares two texts after expansion.
\newcommand{\CheckEqual}[3]{%
	\protected@edef\Check@a{#2}%
	\protected@edef\Check@b{#3}%
	\ifx\Check@a\Check@b
		\typeout{TEST OK: #1}%
	\else
		\typeout{TEST FAILED: #1: expected "\Check@b", got "\Check@a"}%
	\fi
}

% \CheckColour{name}{R,G,B} checks the definition of a colour; the name .
% stands for the colour of the text at the place of the check.
\newcommand{\CheckColour}[2]{%
	\extractcolorspecs{#1}{\Check@a}{\Check@b}%
	\convertcolorspec{\Check@a}{\Check@b}{RGB}{\Check@c}%
	\CheckEqual{colour #1}{\Check@c}{#2}%
}

\makeatother
